\section{Conclusion}
A successful implementation of an autonomous chess playing robot, has been made. The use of computer vision made it able to dynamically find the location of the chessboard on the table, making the player able to move the board or rotate it however desired. This interface also gave the user a choice of whether to be black or white, simply by rotating the chessboard. A validation for every possible move the camera detects was made to further increase the consistency of the move detection. 

While ROS2 was used and provided helpful tools like the simple interface of the UR driver and the GUI in Rviz, it could have been substituted for a simpler integration to reduce overhead. The gripper was made specifically for picking up the used chess pieces, and was made wide to account for any error when locating the chessboard, to account for the offset from the center of the chessboard tiles. However, the gripper's ability to combat any offset of the tile's center, was not enough as the gripper managed to pick up the pieces $97.2 \%$ of the time, in comparison to the desired $98\%$ of the time. The camera in combination with validating each move, resulted in the system correctly detecting the player's move $98.7 \%$ of the time. These 2 percentages result in a $11.33 \%$ chance of completing a game of chess with $50$ moves for both black and white. Further improvements are therefore required to consistently play an autonomous game of chess against a robot.